\documentclass[a4paper,11pt]{article}

\usepackage{revy}
\usepackage[utf8]{inputenc}
\usepackage[T1]{fontenc}
\usepackage[danish]{babel}
\usepackage{graphicx}


\revyname{MatematikRevyen}
\revyyear{2026}
% HUSK AT OPDATERE VERSIONSNUMMER
\version{0.1}
\eta{$5$ minutter}
\status{Færdig}

\title{Jobsamtale på Hilberts Hotel}
\author{Julius Pallesen'24}

\begin{document}
\maketitle

\begin{roles}
\role{M} Matematikkandidat
\role{H} Hotelejer/receptionist Hilbert
\end{roles}

\begin{props}
    \prop{Reception} F.eks. et højt bord med et skilt
    \prop{CV} Curriculum Vitae
\end{props}


\scene{Note: Dette er første sketch i en serie af tre. HH1, HH2 og HH3a kan med minimal omskrivning også sagtens fungere enkeltvis.}
\begin{sketch}
\scene{Scenen er reception på Hilberts Hotel. Der er en receptionist (Hilbert) som har stort facadesmil på og snakker meget hurtigt. Og en jobsøgende matematikkandidat.}

\says{H} Velkommen verdens største hotelkæde, selv kun med ét hotel! Alle værelser er p.t. fyldt op, men det betyder ikke, at vi ikke har plads til DIG! Der er altid plads til en gæst her på Hilberts Hotel! Hvad er navnet?

\says{M} He.. hej jeg hedder M, men... \act{Afbrydes}

\says{H}  Hmm hmm, ah ja her er din booking. Du skal simpelthen bo på værelse 1...

\says{M}  Nej nej du misforstår jeg har ikke booket.

\says{H} Nej det har vi da gjort for dig. Det gør vi for alle. Så længe alle er en tællelig mængde HAHAHA.

\says{M} Nej! Jeg er her fordi jeg har søgt job her og skal til jobsamtale.

\scene{H slipper facadesmilet da det ikke er en kunde.}

\says{H} Ja så følg med mig.

\scene{De sætter sig ved et samtalebord.}

\says{H} Er det dit CV du har der? Må lige se det engang?

\says{M} \act{rækker CV}

\says{H} \act{Skimmer} Hm, ja, ser man det. En kandidatgrad i Matematik jaer. Sådan en har jeg også... OG en bachelorgrad, hmm... Og du har OGSÅ en gymnasial uddannelse, mm... og også bestået folkeskolen. Haha. Du har rigtig været hele møllen igennem hva'?

\says{M} jaeh?

\says{H} Så du er god til matematik, siger du? Hvad er 53.196 divideret med 33?

\says{M} Øhh..

\says{H} \act{Inden M har en chance for at nå at svare} Du er måske ikke så hurtig? Du skal ikke være nervøs. Hvad er 15 i anden gange ti niendedele?

\says{M} $16,5$?

\says{H} Vi går videre... \act{læser videre} Du arbejdet i Rema1000, da du var 16, og gået med aviser da du var 14... Altså det virker ikke lige som om du har så meget erfaring i hotelbranchen?

\says{M} Nej, men...

\says{H} Hvilke kurser har du taget på Matematik

\says{M} Øh, MatIntro...

\says{H} Nej! Kandidatkurser! Når ja det står her: Asgers Logikkursus, Noget mængdelære, APP? \act{Lægger CV'et væk.} Nå, kan du så fortælle mig; hvad kender du egentlig til lige netop dette hotel?

\says{M} Altså.. der er uendelig mange værelser, men det er uendelig mange optagede værelser.

\says{H} Korrekt! Dette hotel har flere værelser OG en større omsætning end alle andre hoteller i verden. Tilsammen! Kan du ikke fortælle lidt mere om dig selv. Hvad ville du sige er din største svaghed?

\says{M} Svagheder? Uha jeg har mange.. Jeg kommer tit for sent. Jeg er asocial. Jeg er blevet postet mere på matgo'nat end nogen anden, fordi jeg altid falder i søvn i timerne. Jeg er afhængig af at spise japansk kridt. Eller.. jeg har gjort det én.. højest to gange.

\says{H} Tilgengæld er du ærlig. Kan du forklarer mig hvordan vi får plads til nye gæster når nu alle værelser til enhver tid er optagede?

\says{M} \act{har ventet på at netop dette spørgsmål skulle komme} Man giver dem værelse 1. Og så beder man personen der var på værelse 1 om at rykke til værelse 2 og så videre. Man beder personen på værelse $n$ om at rykke til værelse $n+1$. Og hvis der kommer en uendelig bus med nummererede sæder, så sender man beboeren på værelse $n$ til $2n-1$, og personen på bussæde $m$ til værelse $2m$.

\says{H} Og hvad gør en uendelig lang bus så, når den kommer til en rundkørsel?

\says{M} Hvad?

\says{H} Anyway.. Her på falderebet, har du så nogen spørgsmål til mig?

\says{M} Er $\aleph_0+1$ kongruent med 0 modulo 4?

\says{H} Hva'? \act{tænker hårdt under M's næste replik}

\says{M} Det er bare fordi jeg \textit{elsker} at spille whist, og tænkte jeg kunne spille med hotelgæsterne i mine pauser. Og jeg vil ikke have nogen føler sig ekskluderet. Og siden der er $\aleph_0$ hotelgæster og jeg er én, skal $\aleph_0+1$ være kongruent med 0 modulo 4 for at det går op.

\says{H} \act{måløs, tænker højt} $\aleph_0$ gange 4, giver $\aleph_0$, så derfor er $\aleph_0$ kongruent med 0 modulo 4, men så er $\aleph_0+1$ jo kongruent med \textit{1} modulo 4, men siden $\aleph_0+1$ jo er \textit{lig} med $\aleph_0$, så må det jo betyde at 1 er lig med 0 modulo 4... Modstrid! \act{Vender et seriøst blik mod M} Hvad er det for et kaos du har skabt unge knægt? Jeg tror vi er færdige for nu.

\says{M} Nå.. okay. Men får jeg så jobbet?

\says{H} \act{Suk} Ja du får jobbet. Men kun fordi der altid er brug for en ekstra hånd her på Hilberts Hotel... Der er dog én betingelse: Du dropper alle ideer om at spille Whist her på matriklen.

\scene{Lys ned.}

\says{AV} Fortsættes...

\end{sketch}

\scene{Note: Hvis denne vælges enkeltvis, ville jeg slette whist-plottet og erstatte det med en punchline.}

\end{document}
